\documentclass[10pt,a4paper]{amsart}
\usepackage[margin=19mm]{geometry}
\usepackage{amsmath,amssymb,mathtools}
\usepackage[scr=rsfs,cal=euler]{mathalfa}
\usepackage{xfrac}
\usepackage[unicode,hidelinks,bookmarks=false]{hyperref}
\newcommand{\ie}{i.e.\ }
\newcommand{\cf}{cf.\ }
\newcommand{\resp}{respectively\ }
\newcommand{\Subsection}{\subsection}
\providecommand{\nobreakdash}{\nobreak\hbox{-}}
\setlength{\emergencystretch}{2em}
\multlinegap0pt
\usepackage{bm}
\usepackage{bbm}
\usepackage{dsfont}
\usepackage{xspace}
\usepackage{cancel}
\usepackage{mathtools}
\usepackage{amsthm}
\makeatletter
\@ifundefined{theorem}{\newtheorem{theorem}{Theorem}}{}
\@ifundefined{assumption}{\newtheorem{assumption}[theorem]{Assumption}}{}
\makeatother
\makeatletter
\expandafter\ifx\csname acknowledgement\endcsname\relax
\newenvironment{acknowledgement}{\par\addvspace{17pt}\small\rmfamily\noindent}{\par\addvspace{6pt}}
\fi
\makeatother
\title[Stable Positive Integral Deferred Correction Methods for Pos]{Stable Positive Integral Deferred Correction Methods for Positive Dynamical Systems}
\author{Pezzella Mario}
\date{Source: arXiv:2606.29523v1 (CC BY 4.0)}
\begin{document}
\maketitle

\noindent\emph{Source and licence.} Stable Positive Integral Deferred Correction Methods for Positive Dynamical Systems. Version arXiv:2606.29523v1 (CC BY 4.0), distributed under the Creative Commons Attribution 4.0 International licence, as recorded in the arXiv licence metadata for this exact version and shown on the abstract page https://arxiv.org/abs/2606.29523v1. Full rights notice: licenses/arxiv-2606.29523-CC-BY-4.0.txt.\par\smallskip

\noindent\textit{Editorial scope.} Counted unit: the Theorem whose source label is \texttt{thm:NonlinContr} in the article (source0.tex lines 962--993, source label \texttt{thm:NonlinContr}), delivered together with the article's own complete original proof of it (source0.tex lines 995--1084, one \texttt{proof} environment of 90 lines). No mathematics is written, repaired or re-derived: statement and proof are the article's own text line for line. Required context retained uncounted: eq:ODE\_system (lines 99--104); eq:L2 (lines 195--205); ass:Solv (lines 406--421); eq:coll\_syst\_Con\_L (lines 409--417); eq:CollocationMap (lines 429--437); eq:LogSystem (lines 884--889); prop:RK\_Log (lines 903--915); eq:RK\_Log (lines 905--910); eq:OSLC (lines 935--944); eq:distance (lines 951--953). Every definition, display and label of those lines is retained unchanged. Omitted from this package and not redistributed: 1--98, 105--194, 206--405, 418--428, 438--883, 890--902, 911--934, 945--950, 954--961, 994--994, 1085--1723. Nothing inside a retained excerpt is omitted or altered: every retained body is delivered exactly as the article's lines are written, so no omission receipt of any kind accompanies this package. Citations: the retained excerpts carry no citation command at all, so no bibliography entry of the article is transcribed and no reference is redistributed. Reference adaptation: none was needed and none was made; every reference inside the delivered unit resolves inside it, so no unresolved reference or citation is produced. Printed slips kept literal: no printed word, symbol, hypothesis or number of the counted statement or of its proof was altered. Layout and class: the article's own class is \texttt{oup-contemporary-spidec} with packages including graphicx, siunitx, amssymb, amsmath, amsthm, amscd, bm, bigints, hyperref, orcidlink; the wrapper is the lane's pinned \texttt{amsart} editorial wrapper at 10pt on A4, which restates the theorem-like environments the retained text needs and adds the article's own macro definitions that the retained text needs (none), with extra wrapper packages (bm, bbm, dsfont, xspace, cancel, mathtools). The delivered numbering is the unit's own. Assets: no retained span contains a figure environment, a graphics-inclusion command, a PDF-page inclusion, a table or a TikZ picture; no image file is copied into this package, the package declares no asset dependency and no image file is redistributed with it. Licence: version arXiv:2606.29523v1 of this article is distributed under the Creative Commons Attribution 4.0 International licence, as recorded in the arXiv licence metadata for this exact version and shown on the abstract page https://arxiv.org/abs/2606.29523v1. A full rights notice is delivered at licenses/arxiv-2606.29523-CC-BY-4.0.txt. Characters: every retained excerpt is pure ASCII, so the delivered engine, XeLaTeX with its default Latin Modern text font, reports no missing character.
\begin{equation}\label{eq:ODE_system}
\begin{cases}
    \bm{y}'(t) = \bm{f}(\bm{y}(t)), \qquad  t \in [0,T],\\
    \bm{y}(0) = \bm{y}_0 \in \mathbb{R}_{+}^{N},
\end{cases}
\end{equation}

\begin{equation}\label{eq:L2}
    \mathcal{L}_i^2(\bm{Y}^{n})=\begin{pmatrix}
        y_i^{n,0}-y_i^{n}\exp\!\left(\displaystyle\int_{t_n}^{t_{n,0}} \mathcal{E}_M^{-1}[\bm{F}^i](s) \, ds \right) \\
        \vdots \\
        y_i^{n,M}-y_i^{n}\exp\!\left(\displaystyle\int_{t_n}^{t_{n,M}} \mathcal{E}_M^{-1}[\bm{F}^i](s) \, ds \right)
    \end{pmatrix}^\top=\begin{pmatrix}
        y_i^{n,0}-y_i^{n}\exp\!\left( h \sum_{j=0}^{M} Q_{0,j} \, \dfrac{f_i(\bm y^{n,j})}{y_i^{n,j}} \right) \\
        \vdots \\
        y_i^{n,M}-y_i^{n}\exp\!\left( h \sum_{j=0}^{M} Q_{M,j}\, \dfrac{f_i(\bm y^{n,j})}{y_i^{n,j}} \right)
    \end{pmatrix}^\top, \quad \quad i=1,\dots,N,
\end{equation}

\begin{assumption}
\label{ass:Solv}
     For each $h>0$ and each initial datum $\bm{y}^n\in\mathbb{R}_{+}^{N}$, the collocation system
    \begin{equation}\label{eq:coll_syst_Con_L}
            \bm{\mathcal{L}}^2(\bm{Y}^n)=\bm{0},   \\
        \qquad \text{where} \qquad \bm{Y}^{n} =
        \big[ \bm{y}^{n,0},\dots,\bm{y}^{n,M} \big]=\begin{bmatrix}
        y_1^{n,0} & \dots & y_1^{n,M} \\
        \vdots & \vdots  & \vdots \\
        y_N^{n,0} & \dots & y_N^{n,M} \\ 
        \end{bmatrix},
    \end{equation}
    with $\bm{y}^n$ entering as a known parameter on the 
    right-hand side of~\eqref{eq:L2}, admits a unique 
    solution $\tilde{\bm{Y}}^n\in\mathbb{R}^{N\times(M+1)}_{>0}$.
\end{assumption}

\begin{equation}\label{eq:CollocationMap}
    \bm{\Psi}^h : \mathbb{R}_{+}^{N} \to \mathbb{R}_{+}^{N},
    \qquad
    \bm{\Psi}^h(\bm{y}^n) =
    \bigl[\Psi_1^h(\bm{y}^n),\dots,\Psi_N^h(\bm{y}^n)\bigr]^\top,
    \qquad
    \Psi_i^h(\bm{y}^n)=\tilde{y}_i^{n,M},
    \quad i=1,\dots,N,
\end{equation}

\begin{equation}\label{eq:LogSystem}
    u^{\prime}_i(t) = G_i(\bm{u}(t)),
    \qquad \text{with} \qquad
    G_i(\bm{u}) = \dfrac{f_i(\operatorname{Exp}(\bm{u}))}{\exp(u_i)},
    \qquad i=1,\dots,N,
\end{equation}

\begin{theorem}\label{prop:RK_Log}
    Under the Assumption~\ref{ass:Solv}, let $\tilde{\bm{Y}}^n=[\tilde{\bm{y}}^{n,0},\dots,\tilde{\bm{y}}^{n,M}]$ denote the unique collocation solution to \eqref{eq:coll_syst_Con_L}. Then, the logarithmic formulation $\tilde{u}_i^{n,m} = \log(\tilde{y}_i^{n,m})$ satisfies the implicit Runge--Kutta system
    \begin{equation}\label{eq:RK_Log}
        \tilde{u}_i^{n,m}=  u_i^n + h\sum_{j=0}^M Q_{m,j}\,G_i\!\left(\tilde{\bm{u}}^{n,j}\right), \quad \begin{array}{l}
             i=1,\dots,N,  \\
             m=0,\dots,M, 
        \end{array} \quad \text{with output} \quad u_i^{n+1} = u_i^n + h\sum_{j=0}^M b_j\, G_i\!\left(\tilde{\bm{u}}^{n,j}\right), \quad b_j := Q_{M,j}.
    \end{equation}
    % with output
    % \begin{equation}\label{eq:RK_Output}
    %     u_i^{n+1} = u_i^n + h\sum_{j=0}^M b_j\, G_i\!\left(\tilde{\bm{u}}^{n,j}\right), \qquad b_j := Q_{M,j}.
    % \end{equation}
\end{theorem}

    \begin{equation}\label{eq:OSLC}
        \langle
          \bm{G}(\bm{u}) - \bm{G}(\bm{v}),\,
          \bm{u} - \bm{v}
        \rangle
        \leq
        \nu\,\|\bm{u}-\bm{v}\|^2,
        \qquad
        \forall\,\bm{u},\bm{v}\in\mathbb{R}^N.
    \end{equation}

\begin{equation}\label{eq:distance}
    \operatorname{d}(\bm{y},\bm{z})= \|\operatorname{Log}(y\oslash\bm{z})\|= \|\operatorname{Log}(\bm{y})-\operatorname{Log}(\bm{z})\|, \qquad \text{where} \qquad \operatorname{Log}(\bm{v})=[\log(v_1),\dots,\log(v_N)]^\top,
\end{equation}

\begin{theorem}[Contractivity of the collocation map]
\label{thm:NonlinContr}
    Let Assumption~\ref{ass:Solv} hold and suppose that 
    the system~\eqref{eq:ODE_system} is logarithmically 
    dissipative with constant $\nu\leq 0$. Assume that 
    the RK scheme~\eqref{eq:RK_Log} is algebraically 
    stable, that is
    \begin{equation}\label{eq:Algebraic_Stability}
        b_j = Q_{M,j} \geq 0
        \quad \text{for each } j=0,\dots,M,
        \qquad \text{and} \qquad
        \bm{\mathcal{M}}
        = \operatorname{diag}(\bm{b})\,Q
          + Q^\top\operatorname{diag}(\bm{b})
          - \bm{b}\bm{b}^\top
        \ \text{is positive semi-definite,}
    \end{equation}
    where $\bm{b}=[b_0,\dots,b_M]^\top$.
    Then, for any $h>0$ and any 
    $\bm{y},\bm{z}\in\mathbb{R}_{+}^{N}$,
    \begin{equation}\label{eq:ContractColloc}
        \operatorname{d}\bigl(
          \bm{\Psi}^h(\bm{y}),\,
          \bm{\Psi}^h(\bm{z})
        \bigr)
        \leq \operatorname{d}(\bm{y},\bm{z}),
    \end{equation}
    where $\operatorname{d}$ denotes the log-Euclidean 
    metric~\eqref{eq:distance}. In particular, 
    the map $\bm{\Psi}^h$ in \eqref{eq:CollocationMap} is non-expansive for $\nu=0$ and 
    strictly contractive for $\nu<0$.
\end{theorem}

\begin{proof}
    By Theorem~\ref{prop:RK_Log}, the map $\bm{\Psi}^h$ 
    acts in log variables as the implicit RK 
    method~\eqref{eq:RK_Log} applied 
    to~\eqref{eq:LogSystem}. Let $\tilde{\bm{u}}^m$ and 
    $\tilde{\bm{v}}^m$ denote the stage values 
    corresponding to initial data 
    $\bm{u}^n=\operatorname{Log}(\bm{y})$ and 
    $\bm{v}^n=\operatorname{Log}(\bm{z})$, and define
    \[
        \bm{\theta} = \bm{u}^n - \bm{v}^n,
        \qquad
        \bm{\gamma}^m = \tilde{\bm{u}}^m 
                        - \tilde{\bm{v}}^m,
        \qquad
        \bm{\kappa}^m = \bm{G}(\tilde{\bm{u}}^m)
                        - \bm{G}(\tilde{\bm{v}}^m),
        \qquad m=0,\dots,M.
    \]
    Note that $\|\bm{\theta}\|
    = \operatorname{d}(\bm{y},\bm{z})$ and 
    $\|\bm{\gamma}^M\|
    = \operatorname{d}(\bm{\Psi}^h(\bm{y}),
                      \bm{\Psi}^h(\bm{z}))$.
    Subtracting the RK systems~\eqref{eq:RK_Log} for 
    the two initial data gives
    \[
        \bm{\gamma}^m 
        = \bm{\theta} 
          + h\sum_{j=0}^M Q_{m,j}\,\bm{\kappa}^j,
        \qquad m=0,\dots,M,
        \qquad\text{and}\qquad
        \bm{\gamma}^M 
        = \bm{\theta} 
          + h\sum_{j=0}^M b_j\,\bm{\kappa}^j.
    \]
    Expanding $\|\bm{\gamma}^M\|^2$ and substituting 
    $\bm{\theta} = \bm{\gamma}^j 
    - h\sum_{\ell=0}^M Q_{j,\ell}\bm{\kappa}^\ell$ yields
    \begin{equation}\label{eq:NormExpansion}
        \begin{split}
        \|\bm{\gamma}^M\|^2
        &= \|\bm{\theta}\|^2
           + 2h\sum_{j=0}^M b_j
             \langle\bm{\theta},\bm{\kappa}^j\rangle
           + h^2\Bigl\|\sum_{j=0}^M 
             b_j\bm{\kappa}^j\Bigr\|^2 \\
        &= \|\bm{\theta}\|^2
           + 2h\sum_{j=0}^M b_j
             \Bigl(\langle\bm{\gamma}^j,\bm{\kappa}^j\rangle
             - h\sum_{\ell=0}^M Q_{j,\ell}
               \langle\bm{\kappa}^\ell,
                       \bm{\kappa}^j\rangle\Bigr)
           + h^2\sum_{j=0}^M\sum_{\ell=0}^M
             b_j b_\ell
             \langle\bm{\kappa}^j,\bm{\kappa}^\ell
             \rangle \\
        &= \|\bm{\theta}\|^2
           + 2h\sum_{j=0}^M b_j
             \langle\bm{\gamma}^j,\bm{\kappa}^j\rangle
           - h^2\sum_{j=0}^M\sum_{\ell=0}^M
             \mathcal{M}_{j,\ell}
             \langle\bm{\kappa}^j,
                     \bm{\kappa}^\ell\rangle,
        \end{split}
    \end{equation}
    where $\mathcal{M}_{j,\ell} 
    = b_j Q_{j,\ell} + b_\ell Q_{\ell,j} - b_j b_\ell$.
    Since $\mathcal{M}\succcurlyeq 0$ by algebraic 
    stability, the last term is non-negative. The 
    one-sided Lipschitz condition~\eqref{eq:OSLC} gives
    \[
        \langle\bm{\gamma}^j,\bm{\kappa}^j\rangle
        = \langle\tilde{\bm{u}}^j - \tilde{\bm{v}}^j,\,
                  \bm{G}(\tilde{\bm{u}}^j)
                  - \bm{G}(\tilde{\bm{v}}^j)\rangle
        \leq \nu\|\bm{\gamma}^j\|^2 \leq 0.
    \]
    Since $\nu\leq 0$ and $b_j\geq 0$, the middle term 
    in~\eqref{eq:NormExpansion} is also non-positive. 
    Therefore $\|\bm{\gamma}^M\|^2 \leq \|\bm{\theta}\|^2$, 
    and taking square roots yields~\eqref{eq:ContractColloc}.
    
    If $\nu<0$, then $\langle\bm{\gamma}^j,\bm{\kappa}^j\rangle < 0$ whenever $\bm{\gamma}^j\neq\bm{0}$. Since $y \neq z$ implies $\theta \neq 0$, and Assumption \ref{ass:Solv} guarantees that the collocation 
solution depends continuously and injectively on 
the initial datum $\bm{u}^n$, we have 
$\gamma^j = \tilde{\bm{u}}^j - \tilde{\bm{v}}^j 
\neq \bm{0}$ for at least one $j\in\{0,\dots,M\}$, 
giving strict inequality.
\end{proof}

\end{document}
